\documentclass[12pt, a4paper]{article}

% ── Packages ──────────────────────────────────────────────────────
\usepackage[utf8]{inputenc}
\usepackage[T1]{fontenc}
\usepackage[margin=2.5cm]{geometry}
\usepackage{graphicx}
\usepackage{hyperref}
\usepackage{listings}
\usepackage{xcolor}
\usepackage{booktabs}
\usepackage{tabularx}
\usepackage{amsmath}
\usepackage{float}
\usepackage{fancyhdr}
\usepackage{titlesec}
\usepackage{parskip}
\usepackage{microtype}
\usepackage{enumitem}
\usepackage{caption}
\usepackage{subcaption}
\usepackage{lmodern}

% ── Hyperref setup ────────────────────────────────────────────────
\hypersetup{
    colorlinks=true,
    linkcolor=black,
    urlcolor=blue,
    citecolor=black
}

% ── Header / footer ───────────────────────────────────────────────
\pagestyle{fancy}
\fancyhf{}
\rhead{AI-Driven IoT Driver Drowsiness Detection}
\lhead{Final Year Project}
\rfoot{Page \thepage}
\renewcommand{\headrulewidth}{0.4pt}

% ── Code listing style ────────────────────────────────────────────
\definecolor{codegray}{rgb}{0.95,0.95,0.95}
\definecolor{codeblue}{rgb}{0.0,0.3,0.7}
\definecolor{codegreen}{rgb}{0.1,0.5,0.1}
\definecolor{codepurple}{rgb}{0.5,0.0,0.5}

\lstdefinestyle{mystyle}{
    backgroundcolor=\color{codegray},
    commentstyle=\color{codegreen},
    keywordstyle=\color{codeblue}\bfseries,
    stringstyle=\color{codepurple},
    basicstyle=\ttfamily\footnotesize,
    breakatwhitespace=false,
    breaklines=true,
    captionpos=b,
    keepspaces=true,
    showspaces=false,
    showstringspaces=false,
    showtabs=false,
    tabsize=2,
    frame=single,
    rulecolor=\color{gray!40}
}
\lstset{style=mystyle}

% ── Title section ─────────────────────────────────────────────────
\title{
    \vspace{-1cm}
    \Huge\textbf{AI-Driven IoT Driver Drowsiness Detection System} \\[0.5cm]
    \large Final Year Project Report \\[0.3cm]
    \normalsize Internet of Things — Department of Computer Engineering
    \vspace{1cm}
}
\author{
    \textbf{Sudenur Tilla} \\
    \small Department of Computer Engineering \\
    \small Istanbul Kultur University \\
    \small Istanbul, Turkiye \\
    \small 2200007172@stu.iku.edu.tr \\
    \and
    \textbf{Hilda Doğa Arslanpençesi} \\
    \small Department of Computer Engineering \\
    \small Istanbul Kultur University \\
    \small Istanbul, Turkiye \\
    \small 2100004850@stu.iku.edu.tr \\
    \and
    \textbf{Busenaz Elik} \\
    \ Department of Computer Engineering \\
    \small Istanbul Kultur University \\
    \small Istanbul, Turkiye \\
    \small 2100005914@stu.iku.edu.tr \\
}
\date{
    \vspace{6cm}
    \normalsize \textbf{Project Advisor} \\[0.2cm]
    \normalsize Mehmet Fatih Yüce \\[0.2cm]
    \normalsize Department of Computer Engineering \\
    \normalsize İstanbul Kültür University \\[0.3cm]
    \today
}

% ══════════════════════════════════════════════════════════════════
\begin{document}
% ══════════════════════════════════════════════════════════════════

\maketitle
\thispagestyle{empty}
\vspace{1cm}

\newpage
% ── Abstract (English) ────────────────────────────────────────────
\begin{abstract}
Driver fatigue is a leading cause of road accidents worldwide. This report presents
an end-to-end IoT system that detects driver drowsiness in real time by monitoring
eye closure, blink rate, head tilt, and hand position through a webcam, and
transmitting the resulting data over a secured UDP channel to a Flask-based backend.
The system employs a two-layer AI approach: a rule-based drowsiness scorer that
produces an explainable 0--100 risk score, and an Isolation Forest model for
unsupervised anomaly detection. Alerts are triggered when either layer signals
danger, accompanied by an audio alarm. A Grafana dashboard provides live
visualization of all metrics. Security is enforced through HMAC-SHA256 packet
signing, nonce-based replay attack prevention, and API key authentication.
Network reliability is evaluated through controlled packet-loss, latency, and
burst-traffic experiments whose results are captured with Wireshark.
\end{abstract}

\newpage

% ── Özet (Türkçe) ─────────────────────────────────────────────────
\renewcommand{\abstractname}{Özet}
\begin{abstract}
Sürücü yorgunluğu, dünya genelinde trafik kazalarının önde gelen nedenlerinden
biridir. Bu rapor, göz kapanma süresi, göz kırpma hızı, baş eğimi ve el
konumunu web kamerası aracılığıyla gerçek zamanlı olarak izleyen ve elde edilen
verileri güvenli bir UDP kanalı üzerinden Flask tabanlı bir arka uca ileten
uçtan uca bir IoT sistemini sunmaktadır. Sistem iki katmanlı bir yapay zeka
yaklaşımı benimsemektedir: açıklanabilir bir 0--100 risk puanı üreten kural
tabanlı uyuşukluk puanlayıcısı ve denetimsiz anomali tespiti için Isolation
Forest modeli. Her iki katmandan herhangi biri tehlike sinyali verdiğinde,
sesli alarm eşliğinde uyarı tetiklenmektedir. Grafana panosu tüm metriklerin
canlı görselleştirmesini sağlamaktadır. Güvenlik; HMAC-SHA256 paket imzalama,
nonce tabanlı tekrar saldırısı önleme ve API anahtarı kimlik doğrulama ile
sağlanmaktadır. Ağ güvenilirliği, Wireshark ile yakalanan paket kaybı, gecikme
ve ani trafik artışı deneyleriyle değerlendirilmektedir.
\end{abstract}

\newpage

\vspace{0.5cm}
\tableofcontents
\newpage

% ══════════════════════════════════════════════════════════════════
\section{Introduction}
% ══════════════════════════════════════════════════════════════════

\subsection{Motivation}

Driver drowsiness is one of the most dangerous and underreported factors in
road traffic accidents. According to the World Health Organization, fatigue
is estimated to be a contributing factor in up to 20\% of all road collisions
worldwide~\cite{who2018}. Unlike alcohol impairment, drowsiness leaves no
measurable trace after an accident, making prevention, rather than
post-incident analysis, the only viable strategy.

Modern commercial driver monitoring systems (Bosch DMS, Seeing Machines,
Mobileye) achieve high accuracy but rely on proprietary hardware, dedicated
infrared cameras, and closed-source software. This makes them inaccessible
for academic research, low-cost deployment, or independent evaluation.

This project addresses the question: \textit{Can a complete, real-time
drowsiness detection pipeline be built using only open-source software and
a standard laptop webcam?}

\subsection{Project Goals}

The system was designed with the following goals in mind:

\begin{enumerate}
    \item \textbf{Real IoT architecture} — sensor data must travel over actual
          UDP network datagrams (not function calls), so that every packet is
          visible in Wireshark and network conditions can be studied.

    \item \textbf{Two-layer AI} — combine a transparent, rule-based scorer
          (explainable to a non-technical audience) with an unsupervised
          machine learning model (Isolation Forest) that catches subtle
          anomalies the rules alone would miss.

    \item \textbf{Dual input modes} — a live webcam module and a synthetic
          simulator must be fully interchangeable, so the system can be
          demonstrated and tested without a camera.

    \item \textbf{Production-grade security} — every UDP packet must be
          cryptographically signed and protected against replay attacks,
          mirroring real IoT deployments where the network cannot be trusted.

    \item \textbf{Live visualization} — a Grafana dashboard must display all
          sensor metrics in real time, updating every 5 seconds.
\end{enumerate}

\subsection{Report Structure}

The remainder of this report is organized as follows.
Section~\ref{sec:arch} presents the overall system architecture and data flow.
Section~\ref{sec:camera} details the camera-based detection module including
the EAR algorithm and head tilt detection.
Section~\ref{sec:simulator} describes the synthetic sensor simulator.
Section~\ref{sec:network} explains the UDP transport layer and bridge.
Section~\ref{sec:api} covers the Flask REST API.
Section~\ref{sec:ai} presents the two-layer AI detection approach.
Section~\ref{sec:security} describes the security mechanisms.
Section~\ref{sec:experiments} reports the network impairment experiments.
Section~\ref{sec:grafana} covers the Grafana dashboard.
Section~\ref{sec:alerts} explains the alert system.
Section~\ref{sec:results} evaluates overall system performance.
Section~\ref{sec:discussion} discusses design decisions and limitations.
Section~\ref{sec:conclusion} concludes with future work directions.

% ══════════════════════════════════════════════════════════════════
\section{System Architecture}
\label{sec:arch}
% ══════════════════════════════════════════════════════════════════

\subsection{Overview}

The system follows a layered IoT pipeline with four distinct tiers:

\begin{enumerate}
    \item \textbf{Sensing} — webcam or synthetic simulator produces sensor readings.
    \item \textbf{Transport} — UDP datagrams carry HMAC-signed JSON payloads.
    \item \textbf{Processing} — Flask API scores, analyses, and stores readings.
    \item \textbf{Presentation} — Grafana renders live time-series dashboards.
\end{enumerate}

\subsection{Data Flow}

Figure~\ref{fig:architecture} illustrates the complete data flow of the system,
from sensor input to live dashboard visualization.

\begin{figure}[H]
    \centering
    \includegraphics[width=\textwidth]{architecture.png}
    \caption{System architecture and data flow}
    \label{fig:architecture}
\end{figure}

\subsection{Component Summary}

\begin{table}[H]
\centering
\caption{System components and their roles}
\begin{tabularx}{\textwidth}{llX}
\toprule
\textbf{File} & \textbf{Language} & \textbf{Role} \\
\midrule
\texttt{camera/camera\_detector.py} & Python & Dlib EAR webcam detector \\
\texttt{sensor/simulator.py}        & Python & Synthetic 3-state UDP sensor \\
\texttt{network/udp\_bridge.py}     & Python & UDP listener, HMAC verifier, HTTP forwarder \\
\texttt{api/app.py}                 & Python & Flask REST API, scorer, AI runner \\
\texttt{ai/detector.py}             & Python & Isolation Forest training and prediction \\
\texttt{alerts/alert\_manager.py}   & Python & Alert evaluation, storage, audio alarm \\
\texttt{security/crypto.py}         & Python & HMAC-SHA256 sign and verify \\
\texttt{security/replay\_guard.py}  & Python & Nonce-based replay attack prevention \\
\texttt{network/experiments.py}     & Python & Packet loss, latency, burst experiments \\
\bottomrule
\end{tabularx}
\end{table}

% ══════════════════════════════════════════════════════════════════
\section{Camera-Based Detection Module}
\label{sec:camera}
% ══════════════════════════════════════════════════════════════════

\subsection{Eye Aspect Ratio (EAR)}

The camera module uses the Dlib 68-point facial landmark predictor to locate the
six landmark points around each eye. The Eye Aspect Ratio is defined as:

\begin{equation}
    \text{EAR} = \frac{\|p_2 - p_6\| + \|p_3 - p_5\|}{2\,\|p_1 - p_4\|}
\end{equation}

where $p_1 \ldots p_6$ are the six eye landmark points in clockwise order.
When the eye is open, EAR is approximately 0.20--0.35. When the eye closes,
EAR drops toward zero. A threshold of $\text{EAR} < 0.15$ is used to classify
the eye as closed; this value is calibrated per user using \texttt{calibrate\_ear.py}.

\subsection{Calibration Procedure}

Before first use, the user runs a two-phase calibration:

\begin{enumerate}
    \item \textbf{Phase 1 (5 s):} Keep eyes open $\rightarrow$ measure average open EAR.
    \item \textbf{Phase 2 (5 s):} Keep eyes closed $\rightarrow$ measure average closed EAR.
\end{enumerate}

The threshold is set at:
\begin{equation}
    \text{threshold} = \text{EAR}_\text{closed} + 0.35 \times (\text{EAR}_\text{open} - \text{EAR}_\text{closed})
\end{equation}

For the primary test subject: $\text{EAR}_\text{open} = 0.218$,
$\text{EAR}_\text{closed} = 0.114$, giving $\text{threshold} = 0.150$.

\subsection{Detected Signals}

\begin{table}[H]
\centering
\caption{Camera-detected drowsiness signals}
\begin{tabularx}{\textwidth}{lXll}
\toprule
\textbf{Signal} & \textbf{Method} & \textbf{Normal} & \textbf{Drowsy} \\
\midrule
Eye closure duration & EAR below threshold & $<0.4$ s & $>2.0$ s \\
Blink rate           & 60 s rolling window & 15--20 /min & $<6$ /min \\
Head tilt (lateral)  & Eye-corner angle    & $0$--$15°$ & $>25°$ for $>4$ s \\
Head tilt (forward)  & Nose--chin ratio    & normal dist. & ratio drop \\
Both hands raised    & Skin contour area   & 0 contours UP & 2 contours UP $>1.5$ s \\
\bottomrule
\end{tabularx}
\end{table}

\subsection{Head Tilt Detection}

Two tilt axes are monitored simultaneously:

\textbf{Lateral tilt} is computed from the angle formed by the left and right
outer eye corners (landmarks 36 and 45):
\begin{equation}
    \theta_\text{lateral} = \left|\arctan\!\left(\frac{y_R - y_L}{x_R - x_L}\right)\right|
\end{equation}

\textbf{Forward tilt} is detected by measuring the normalised vertical distance
between the nose tip (landmark 30) and chin (landmark 8). As the head drops
forward, this distance decreases:
\begin{equation}
    \theta_\text{forward} = \max\!\left(0,\; \frac{0.13 - d_\text{norm}}{0.13} \times 45\right)
\end{equation}

The larger of the two angles is reported as \texttt{head\_tilt\_angle}.
A duration timer tracks how long the head has been tilted beyond $15°$.

\subsection{UDP Packet Format}

Every second, the camera module transmits:

\begin{lstlisting}[language=Python, caption={Camera UDP payload}]
{
  "driver_id":            "driver_camera_001",
  "state":                "alert",          # alert | transitioning | drowsy
  "blink_rate":           17.2,
  "eye_closure_duration": 0.142,
  "head_tilt_angle":      3.1,
  "head_tilt_duration":   0.0,
  "reaction_delay":       250.0,
  "source":               "camera",
  "nonce":                "a3f2c1d4-...",
  "sent_at":              "2026-05-18T14:33:54+00:00"
}
\end{lstlisting}

% ══════════════════════════════════════════════════════════════════
\section{Synthetic Sensor Simulator}
\label{sec:simulator}
% ══════════════════════════════════════════════════════════════════

The simulator (\texttt{sensor/simulator.py}) generates realistic sensor readings
without requiring a webcam. It cycles through three states:

\begin{table}[H]
\centering
\caption{Simulator state machine}
\begin{tabularx}{\textwidth}{lXl}
\toprule
\textbf{State} & \textbf{Sensor values} & \textbf{Duration} \\
\midrule
\texttt{alert}         & blink 15--20/min, eye $<0.2$ s, tilt $<10°$ & 10 s \\
\texttt{transitioning} & blink 8--15/min, eye 0.2--0.5 s, tilt 10--20° & 8 s \\
\texttt{drowsy}        & blink $<8$/min, eye $>0.5$ s, tilt $>20°$, react $>400$ ms & 12 s \\
\bottomrule
\end{tabularx}
\end{table}

Gaussian noise is added to each value to simulate biological variance. The output
JSON structure is identical to the camera module, so the downstream pipeline
requires no modification to switch between sources.

% ══════════════════════════════════════════════════════════════════
\section{Network Transport Layer}
\label{sec:network}
% ══════════════════════════════════════════════════════════════════

\subsection{Why UDP?}

UDP was chosen over TCP for the following reasons:

\begin{itemize}
    \item \textbf{Low latency:} No connection setup, no acknowledgement round-trips.
    \item \textbf{IoT realism:} Real vehicle telematics systems (CAN bus gateways,
          V2X) use UDP or UDP-like protocols.
    \item \textbf{Observability:} Every datagram is individually visible in Wireshark,
          enabling meaningful packet-loss and latency experiments.
    \item \textbf{Experiment validity:} TCP's automatic retransmission would mask
          packet loss; UDP exposes it directly.
\end{itemize}

\subsection{UDP Bridge}

\texttt{network/udp\_bridge.py} listens on port 9999, performs security checks,
measures per-packet latency, and forwards the payload to the Flask API:

\begin{enumerate}
    \item Receive raw UDP datagram.
    \item Decode JSON envelope \texttt{\{payload, sig\}}.
    \item Verify HMAC-SHA256 signature.
    \item Check nonce has not been seen before (replay guard).
    \item Record timestamp, compute latency.
    \item Forward payload via HTTP POST to \texttt{localhost:5000/sensor-data}.
    \item Log result: packet size, latency, Flask response code, risk level.
\end{enumerate}

\subsection{Observed Latency}

Under normal local conditions the bridge achieves sub-millisecond forwarding
latency (average $\approx 0.4$ ms). This is well within the 1-second sampling
interval, introducing no perceptible delay to the alert pipeline.

% ══════════════════════════════════════════════════════════════════
\section{Backend API}
\label{sec:api}
% ══════════════════════════════════════════════════════════════════

\subsection{Flask Endpoints}

\begin{table}[H]
\centering
\caption{REST API endpoints}
\begin{tabularx}{\textwidth}{llX}
\toprule
\textbf{Method} & \textbf{Endpoint} & \textbf{Description} \\
\midrule
GET  & \texttt{/health}             & Status and total reading count \\
POST & \texttt{/sensor-data}        & Receive and process one reading \\
GET  & \texttt{/readings}           & All stored readings \\
GET  & \texttt{/readings?limit=N}   & Last $N$ readings \\
GET  & \texttt{/readings/summary}   & Average score, risk distribution \\
GET  & \texttt{/ai/analyze}         & AI analysis on last 20 readings \\
GET  & \texttt{/alerts}             & All triggered alerts \\
GET  & \texttt{/alerts/latest}      & Most recent alert \\
POST & \texttt{/query}              & Grafana time-series endpoint \\
POST & \texttt{/metrics}            & Grafana metric list \\
\bottomrule
\end{tabularx}
\end{table}

\subsection{Request Processing Pipeline}

On each \texttt{POST /sensor-data}:

\begin{enumerate}
    \item Input validation with Pydantic (field types, value ranges).
    \item Timestamp assignment (UTC ISO-8601).
    \item Rule-based drowsiness score computation.
    \item Risk level classification (alert / warning / danger).
    \item Isolation Forest prediction (if $\geq 10$ readings available).
    \item Alert evaluation and optional alert creation.
    \item Append reading to \texttt{data/readings.json}.
    \item Return JSON response with score, risk level, and AI flag.
\end{enumerate}
\newpage
% ══════════════════════════════════════════════════════════════════
\section{AI Detection — Two Layers}
\label{sec:ai}
% ══════════════════════════════════════════════════════════════════

\subsection{Layer 1: Rule-Based Drowsiness Scorer}

Each sensor field is mapped to a normalised danger score $s \in [0, 1]$
using a piecewise linear function. The composite score is:

\begin{equation}
    S = 100 \times \bigl(
        0.20\,s_\text{blink}
      + 0.24\,s_\text{eye}
      + 0.20\,s_\text{tilt}
      + 0.16\,s_\text{react}
      + 0.20\,s_\text{phone}
    \bigr)
\end{equation}

The scoring functions are:

\begin{align}
    s_\text{blink}  &= \min\!\left(1,\; \frac{\max(0,\;15 - r_b)}{10}\right) \\
    s_\text{eye}    &= \min\!\left(1,\; \frac{\max(0,\; d - 0.40)}{2.10}\right) \\
    s_\text{tilt}   &= \min\!\left(1,\; \frac{\max(0,\; \theta - 15)}{30}\right) \\
    s_\text{react}  &= \min\!\left(1,\; \frac{\max(0,\; t - 300)}{600}\right)
\end{align}

where $r_b$ is blink rate (blinks/min), $d$ is eye closure duration (s),
$\theta$ is head tilt angle (degrees), and $t$ is reaction delay (ms).

Risk levels are assigned as:

\begin{table}[H]
\centering
\caption{Risk level thresholds}
\begin{tabular}{ccc}
\toprule
\textbf{Score range} & \textbf{Risk level} & \textbf{Indicator} \\
\midrule
0--39  & Alert   & \textcolor{green}{$\bullet$} \\
40--69 & Warning & \textcolor{orange}{$\bullet$} \\
70--100 & Danger & \textcolor{red}{$\bullet$} \\
\bottomrule
\end{tabular}
\end{table}

\subsection{Layer 2: Isolation Forest (Unsupervised Anomaly Detection)}

Isolation Forest~\cite{liu2008isolation} is an ensemble method that isolates
observations by randomly partitioning the feature space. Anomalous points
(drowsy readings) require fewer partitions to isolate than normal points,
resulting in a lower anomaly score.

The model is configured as:

\begin{lstlisting}[language=Python, caption={Isolation Forest configuration}]
IsolationForest(
    n_estimators=100,
    contamination=0.15,   # expect ~15% anomalous readings
    random_state=42
)
\end{lstlisting}

Features used: \texttt{blink\_rate}, \texttt{eye\_closure\_duration},
\texttt{head\_tilt\_angle}, \texttt{reaction\_delay}.

The model is retrained on all available readings with every new incoming packet.
No labels are required — the model learns the distribution of normal readings
and flags deviations automatically.

\subsection{Alert Trigger Conditions}

An alert fires when \textbf{either} condition is met:

\begin{itemize}
    \item Drowsiness score $S \geq 65$
    \item Isolation Forest anomaly score $< -0.08$
\end{itemize}

Using both conditions in parallel reduces false negatives: the rule-based
layer catches clear threshold violations; the AI layer catches subtle
multi-dimensional anomalies that individual thresholds miss.

\subsection{AI Evaluation Results}

After training on 1,847 readings (615 drowsy, 616 alert, 616 transitioning):

\begin{table}[H]
\centering
\caption{Isolation Forest performance on held-out drowsy readings}
\begin{tabular}{lcc}
\toprule
\textbf{Metric} & \textbf{Value} \\
\midrule
Total readings & 1,847 \\
Drowsy readings & 615 \\
Anomaly detection rate (drowsy) & $\approx 65\%$ \\
False positive rate (alert readings) & $<5\%$ \\
\bottomrule
\end{tabular}
\end{table}

The 65\% detection rate reflects the \texttt{contamination=0.15} setting:
only the most extreme 15\% of readings are classified as anomalies. Higher
contamination values improve recall at the cost of more false positives.

\newpage
% ══════════════════════════════════════════════════════════════════
\section{Security Layer}
\label{sec:security}
% ══════════════════════════════════════════════════════════════════

\subsection{Threat Model}

The following threats are addressed:

\begin{itemize}
    \item \textbf{Payload tampering:} An attacker on the local network modifies
          sensor values to suppress or fabricate alerts.
    \item \textbf{Replay attacks:} A captured packet is retransmitted to flood
          the API or replay a clean reading during a drowsy episode.
    \item \textbf{Unauthorised API access:} An unauthenticated client reads or
          injects readings via the REST endpoints.
\end{itemize}

\subsection{HMAC-SHA256 Packet Signing}

Every outgoing packet is wrapped in a signed envelope:

\begin{lstlisting}[language=Python, caption={HMAC envelope structure}]
{
    "payload": { ...sensor fields... },
    "sig":     "<hex HMAC-SHA256>"
}
\end{lstlisting}

The signature is computed over the canonical JSON serialisation of the payload
(keys sorted, no whitespace). The bridge rejects any packet whose signature
does not match. The shared secret (\texttt{IOT\_HMAC\_SECRET}) is stored in
\texttt{.env} and never transmitted.

\subsection{Nonce-Based Replay Prevention}

Each payload includes a UUID4 nonce and a timestamp (\texttt{sent\_at}).
The bridge maintains an in-memory set of seen nonces. Any packet whose nonce
has already been processed is rejected with a log message. This prevents
replay attacks even if a valid signed packet is captured and retransmitted.

\subsection{API Key Authentication}

All Flask endpoints (except \texttt{/health} and \texttt{/metrics}) require
an \texttt{X-API-Key} header matching \texttt{IOT\_API\_KEY} from \texttt{.env}.
The \texttt{@require\_api\_key} decorator is applied at the route level.
Grafana's datasource is configured with the same header.

% ══════════════════════════════════════════════════════════════════
\section{Network Experiments}
\label{sec:experiments}
% ══════════════════════════════════════════════════════════════════

\subsection{Experimental Setup}

Three network conditions were simulated using \texttt{network/experiments.py},
which programmatically introduces impairments at the sender side. Wireshark
was used to capture and verify the traffic on \texttt{udp.port == 9999}.

\subsection{Experiment 1: Packet Loss}

Packets were dropped at the application layer with probabilities from 0\% to 50\%.

\begin{table}[H]
\centering
\caption{Packet loss experiment results}
\begin{tabular}{cccc}
\toprule
\textbf{Loss rate} & \textbf{Sent} & \textbf{Received} & \textbf{Delivery rate} \\
\midrule
0\%  & 20 & 20 & 100\% \\
10\% & 20 & 18 & 90\% \\
25\% & 20 & 15 & 75\% \\
50\% & 20 & 10 & 50\% \\
\bottomrule
\end{tabular}
\end{table}

\textbf{Finding:} At 25\% loss, alert latency increases by approximately 2--3 seconds
as consecutive readings may be missed. At 50\% loss, drowsy episodes of less
than 6 seconds may go undetected.

\subsection{Experiment 2: Injected Latency}

Artificial delays from 0 ms to 500 ms were added before each UDP send.

\begin{table}[H]
\centering
\caption{Latency experiment results}
\begin{tabular}{ccc}
\toprule
\textbf{Injected delay} & \textbf{Bridge latency (avg)} & \textbf{Alert impact} \\
\midrule
0 ms   & 0.4 ms & None \\
100 ms & 100.4 ms & Negligible \\
250 ms & 250.4 ms & Minor ($<1$ s added) \\
500 ms & 500.4 ms & Noticeable delay \\
\bottomrule
\end{tabular}
\end{table}

\textbf{Finding:} Since the sampling interval is 1 second, delays below 500 ms
have no practical impact on alert timing. The system tolerates up to 900 ms
of injected delay without missing a sampling window.

\subsection{Experiment 3: Burst Traffic}

Burst packets (1--50 packets sent simultaneously) were injected to test
buffer handling.

\begin{table}[H]
\centering
\caption{Burst traffic experiment results}
\begin{tabular}{ccc}
\toprule
\textbf{Burst size} & \textbf{Processing time} & \textbf{Packets lost} \\
\midrule
1   & $<1$ ms  & 0 \\
10  & 8 ms     & 0 \\
25  & 21 ms    & 0 \\
50  & 45 ms    & 0 \\
\bottomrule
\end{tabular}
\end{table}

\textbf{Finding:} The bridge processes all burst packets without loss. The
Flask API queues requests during bursts; all readings are eventually stored.

% ══════════════════════════════════════════════════════════════════
\section{Grafana Dashboard}
\label{sec:grafana}
% ══════════════════════════════════════════════════════════════════

The Grafana dashboard provides live monitoring through five time-series panels,
each polling the Flask \texttt{/query} endpoint every 5 seconds via the
\texttt{simpod-json-datasource} plugin.

\begin{table}[H]
\centering
\caption{Grafana dashboard panels}
\begin{tabularx}{\textwidth}{lXl}
\toprule
\textbf{Panel} & \textbf{What it shows} & \textbf{Expected range} \\
\midrule
Drowsiness Score      & Composite 0--100 risk score & 0--40 (alert) \\
Blink Rate            & Blinks per minute (60 s window) & 15--20 \\
Eye Closure Duration  & Eye closed time per event & 0.1--0.4 s \\
Head Tilt Angle       & Degrees from upright & 0--15° \\
Reaction Delay        & Response latency & 150--300 ms \\
\bottomrule
\end{tabularx}
\end{table}

The dashboard is configured with \texttt{Last 1 hour} time range and 5-second
auto-refresh. During a drowsy episode, spikes are immediately visible in the
Drowsiness Score, Eye Closure Duration, and Reaction Delay panels simultaneously.

% ══════════════════════════════════════════════════════════════════
\section{Alert System}
\label{sec:alerts}
% ══════════════════════════════════════════════════════════════════

When an alert condition is met, the system:

\begin{enumerate}
    \item Prints a formatted alert block to the terminal.
    \item Appends a structured record to \texttt{data/alerts.json}.
    \item Plays an audio alarm via \texttt{alerts/sound\_alert.py}.
\end{enumerate}

Each alert record contains:

\begin{lstlisting}[language=Python, caption={Alert record structure}]
{
    "alert_id":         421,
    "timestamp":        "2026-05-18T13:30:46+00:00",
    "driver_id":        "driver_camera_001",
    "drowsiness_score": 67.3,
    "risk_level":       "warning",
    "ai_anomaly":       true,
    "ai_score":         -0.0890,
    "reasons":          ["AI anomaly detected (score: -0.0890)"],
    "sensor_snapshot":  {
        "blink_rate":           8.2,
        "eye_closure_duration": 2.15,
        "head_tilt_angle":      4.3,
        "reaction_delay":       250.0
    }
}
\end{lstlisting}
\newpage

% ══════════════════════════════════════════════════════════════════
\section{Results and Evaluation}
\label{sec:results}
% ══════════════════════════════════════════════════════════════════

\subsection{End-to-End Latency}

The total latency from a drowsy event occurring to an alert being triggered
was measured across 50 trials:

\begin{table}[H]
\centering
\caption{End-to-end alert latency}
\begin{tabular}{lc}
\toprule
\textbf{Stage} & \textbf{Latency} \\
\midrule
Camera frame capture to UDP send & $<33$ ms (30 fps) \\
UDP send to bridge receive & $<1$ ms (local) \\
Bridge to Flask API & $0.4$ ms (avg) \\
Flask processing (score + AI) & $<50$ ms \\
\midrule
\textbf{Total (camera to alert)} & $\mathbf{<500}$ \textbf{ms} \\
\bottomrule
\end{tabular}
\end{table}

\subsection{Detection Performance}

\begin{table}[H]
\centering
\caption{Detection accuracy by input mode}
\begin{tabularx}{\textwidth}{Xcc}
\toprule
\textbf{Condition} & \textbf{Simulator} & \textbf{Camera} \\
\midrule
Eyes closed $>2$ s detected       & 100\% & $\approx 95\%$ \\
Head tilt $>4$ s detected         & 100\% & $\approx 90\%$ \\
Normal state correctly classified & 100\% & $\approx 92\%$ \\
AI anomaly on drowsy readings     & $\approx 65\%$ & $\approx 60\%$ \\
\bottomrule
\end{tabularx}
\end{table}

Camera detection rates are lower than simulator rates due to real-world
factors: varying lighting, face occlusion, and head pose variation. The
EAR-based approach significantly outperformed the initial Haar Cascade
implementation, which produced unstable blink rates of 50--60 blinks/min.

\subsection{Security Verification}

\begin{table}[H]
\centering
\caption{Security mechanism verification}
\begin{tabular}{lcc}
\toprule
\textbf{Attack} & \textbf{Attempted} & \textbf{Blocked} \\
\midrule
Unsigned packet         & 20 & 20 (100\%) \\
Tampered payload        & 20 & 20 (100\%) \\
Replay (duplicate nonce) & 20 & 20 (100\%) \\
Missing API key         & 20 & 20 (100\%) \\
\bottomrule
\end{tabular}
\end{table}
\newpage
% ══════════════════════════════════════════════════════════════════
\section{Discussion}
\label{sec:discussion}
% ══════════════════════════════════════════════════════════════════

\subsection{Design Decisions}

\textbf{Dlib over Haar Cascade:}
The initial implementation used OpenCV Haar Cascade classifiers for eye detection.
This approach produced erratic blink counts (50--60 blinks/min at rest) because
the cascade frequently loses eye detection between frames. Switching to Dlib's
EAR algorithm eliminated this problem — EAR is computed directly from landmark
coordinates and is robust to brief detection gaps.

\textbf{Isolation Forest over supervised learning:}
A supervised classifier would require a large, labelled dataset of real drowsy
driving footage — impractical for a single-person project. Isolation Forest
requires no labels and trains on whatever data is available, making it ideal
for a system that collects its own training data in real time.

\textbf{UDP over TCP:}
While TCP would guarantee delivery, it would also hide the packet loss that
is a defining characteristic of real IoT deployments. UDP makes the network
behaviour explicit and measurable, which is essential for the network
experiments section of this project.

\subsection{Limitations}

\begin{itemize}
    \item \textbf{Skin-colour hand detection} is sensitive to lighting conditions
          and misclassifies background regions. MediaPipe Hands would provide
          much higher reliability but does not support Windows with Python 3.12.
    \item \textbf{Single driver} — the EAR threshold is calibrated per user.
          A production system would require automatic per-session calibration.
    \item \textbf{Reaction delay} is a fixed baseline (250 ms) rather than a
          measured response to a stimulus. A real measurement would require
          a physical button or screen interaction.
    \item \textbf{Isolation Forest retraining} on every incoming packet is
          computationally wasteful. A production system would retrain
          periodically or use an online learning algorithm.
\end{itemize}
\newpage
% ══════════════════════════════════════════════════════════════════
\section{Conclusion}
\label{sec:conclusion}
% ══════════════════════════════════════════════════════════════════

This project demonstrates a complete, production-inspired IoT pipeline for
driver drowsiness detection. The key contributions are:

\begin{enumerate}
    \item A \textbf{dual-input} architecture where simulator and camera modes
          are fully interchangeable without any downstream code changes.
    \item A \textbf{two-layer AI} system combining transparent rule-based
          scoring with unsupervised anomaly detection.
    \item A \textbf{full security stack} — HMAC signing, nonce-based replay
          prevention, and API key authentication — applied to an IoT UDP pipeline.
    \item \textbf{Quantified network experiments} covering packet loss,
          latency, and burst traffic with Wireshark validation.
    \item A \textbf{calibrated EAR-based} camera module using Dlib's 68-point
          landmark model, achieving sub-500 ms alert latency.
\end{enumerate}

Future work could address the identified limitations by integrating a
cross-platform hand landmark model, implementing per-session EAR auto-calibration,
and replacing the fixed reaction delay with a measured stimulus-response protocol.
Deployment on a Raspberry Pi with a dedicated camera module would also
validate the system under genuine embedded IoT constraints.
\newpage
% ══════════════════════════════════════════════════════════════════
% References
% ══════════════════════════════════════════════════════════════════
\newpage
\begin{thebibliography}{9}

\bibitem{liu2008isolation}
F.~T. Liu, K.~M. Ting, and Z.-H. Zhou,
``Isolation Forest,''
in \textit{Proc. 8th IEEE Int. Conf. Data Mining (ICDM)},
pp.~413--422, 2008.

\bibitem{soukupova2016eye}
T.~Soukupova and J.~Cech,
``Real-Time Eye Blink Detection using Facial Landmarks,''
in \textit{21st Computer Vision Winter Workshop}, 2016.

\bibitem{dlib}
D.~E. King,
``Dlib-ml: A Machine Learning Toolkit,''
\textit{Journal of Machine Learning Research}, vol.~10,
pp.~1755--1758, 2009.

\bibitem{who2018}
World Health Organization,
``Global Status Report on Road Safety 2018,''
Geneva: WHO, 2018.

\bibitem{sklearn}
F.~Pedregosa et al.,
``Scikit-learn: Machine Learning in Python,''
\textit{Journal of Machine Learning Research}, vol.~12,
pp.~2825--2830, 2011.

\end{thebibliography}

\end{document}
